A la fin de ce chapitre, je saurai :
\begin{competences}
  \point Exprimer les conditions théoriques (gain et fréquence) d’auto-oscillation sinusoïdale d’un système linéaire bouclé.
  \point Analyser sur l’équation différentielle l’inégalité que doit vérifier le gain de l’amplificateur afin d’assurer le démarrage des oscillations.
  \point Interpréter le rôle des non-linéarités dans la stabilisation de l’amplitude des oscillations.
  \point Décrire les différentes séquences de fonctionnement.
  \point Exprimer les conditions de basculement.
  \point Déterminer l’expression de la période d’oscillation.
  \point \faLaptop À l’aide d’un langage de programmation, simuler l’évolution temporelle d’un signal généré par un oscillateur.
  \point \faSignLanguage Mettre en œuvre un oscillateur quasi-sinusoïdal et analyser les spectres des signaux générés.
\end{competences}
